Large language models must represent millions of compound concepts---``washing machine,'' ``hot dog,'' ``red herring''---yet their residual streams have limited dimensionality.
How do these multi-token concepts fit?
We investigate whether compound nouns are stored as dedicated directions in the residual stream or are constructed compositionally from their parts.
Through three complementary experiments on \gpttwo (124M parameters), we examine 21 compound concepts spanning the compositionality spectrum from transparent (``kitchen chair'') to idiomatic (``red herring'').
We find that compound representations are highly similar to their head nouns in the residual stream (mean cosine similarity 0.954), yet sparse autoencoder (SAE) decomposition reveals that 56.3\% of features active for compounds are unique to the compound context---invisible to cosine similarity.
The degree of unique representation scales with idiomaticity: ``red herring'' activates 82.1\% unique features versus 46.9\% for ``kitchen chair'' (Mann-Whitney $p = 0.0005$, Cohen's $d = 2.54$).
Most strikingly, next-token prediction serves as the dominant storage mechanism: after ``The washing,'' the model assigns 47.1\% probability to ``machine'' (a $4{,}221\times$ increase over baseline), effectively constructing the compound concept through sequential prediction rather than holistic representation.
These findings reveal that LLMs manage the ``too many concepts, too few dimensions'' problem by constructing multi-token concepts sequentially, with implications for interpretability, concept editing, and SAE-based analysis.
